\documentclass[11pt]{amsart}
\usepackage[T1]{fontenc}
\usepackage{lmodern,amsmath,amssymb,mathtools,microtype,booktabs,needspace}
\usepackage[margin=1.15in]{geometry}
\usepackage[hidelinks]{hyperref}
\newtheorem{theorem}{Theorem}[section]
\newtheorem{proposition}[theorem]{Proposition}
\newtheorem{lemma}[theorem]{Lemma}
\newtheorem{corollary}[theorem]{Corollary}
\theoremstyle{definition}\newtheorem{definition}[theorem]{Definition}
\theoremstyle{remark}\newtheorem{remark}[theorem]{Remark}
\newcommand{\R}{\mathbb R}\newcommand{\Z}{\mathbb Z}
\newcommand{\Prob}{\mathbb P}\newcommand{\E}{\mathbb E}
\newcommand{\dd}{\,\mathrm d}\newcommand{\one}{\mathbf 1}
\DeclareMathOperator{\area}{area}\DeclareMathOperator{\covol}{covol}
\DeclareMathOperator{\dist}{dist}\DeclareMathOperator{\conv}{conv}
\DeclareMathOperator{\supp}{supp}
\numberwithin{equation}{section}
\title[Reference-free certification]{Reference-free certification from intrinsic boundary laws}
\author{Qian Qi}\date{October 3, 2026}
\subjclass[2020]{37D50, 52A27, 53A04, 62G05, 62L10}
\keywords{Dispersing billiards, intrinsic boundary laws, first-impact sampling,
finite smoothness, period reconstruction, fixed aperture, sequential certification}
\hypersetup{pdftitle={Reference-free certification from intrinsic boundary laws},pdfauthor={Qian Qi}}
\begin{document}
\raggedbottom
\begin{abstract}
We study whole-table recovery for periodic dispersing billiards from
resettable spatial launches with localized collision positions. A bound
on a lattice presentation permits a fixed-aperture experiment: the
position of the first impact before one fixed short horizon determines
the complete periodic obstacle union and its full translation group.
On a translation-separated $C^{6,\beta}$ class, neither quantitative
asymmetry nor a clear-return network is required. Simultaneous first-hit
coverage, complete-body recognition in the common laboratory frame and
finite rational recovery give $C^2$ geometric error $O(\nu)$ with
$O(\nu^{-3/4}\log(C/(\nu\delta)))$ attempted launches at confidence
$1-\delta$. The aperture is independent of accuracy; the free area is
recovered geometrically, without endpoint histograms or an unknown-cell
normalization. The lattice bound, shape separation, curvature and physical
separation are numerical priors. The sensor still requires controlled
independent launches and calibrated impact-position error; it is not
passive orbit data. We also retain the record-local intrinsic two-density
inverse, completion defect, protected-bridge recognition and anytime
revalidation theorem, with a corrected endpoint-excluding clearance test.
All launch counts include solid starts and no-impact outcomes; precision,
geometric arithmetic and apparatus work are separate. Earlier moment,
analytic and smooth relative-law results remain supplied in full.
\end{abstract}
\maketitle
\section{A fixed-aperture first-impact inverse}
\label{sec:aperture}
The resettable apparatus permits a simpler global observation than a list
of protected return laws. A numerical bound on a lattice presentation
places a translate of every obstacle type, and a generating pair of
translations of any one type, in a fixed laboratory square. Short
first-impact launches in that square therefore see a complete periodic
motif. This observation removes the need to normalize a return law by an
unknown cell area. It also removes the matching of independently oriented
contact frames: all impacts already lie in one laboratory frame.

The resulting theorem concerns the same periodic reflecting bodies and
the same collision apparatus. Its datum is a single first-impact
\emph{position law}, not a scalar collision count, a passive orbit, or an
exterior travelling-time spectrum. Its aperture is fixed by the numerical
prior rather than increased with the target accuracy. The local-return
route in the subsequent sections is retained in full; it gives a different
completion certificate when the acquired clouds have not already covered
a periodic motif.

\subsection{A weaker geometric prior and one position law}
\begin{definition}[Translation-separated periodic class]
\label{def:aperture-prior}
Fix numerical bounds
\[
 r_0,D_0,L_0,M_6,\kappa_+<\infty,\qquad
 d_0,\kappa_-,V_0,\sigma>0,\qquad 0<\beta\le1.
\]
Consider nonempty periodic unions of disjoint strictly convex bodies
\[
 \mathcal O=\bigcup_{i=1}^r(C_i+\Lambda),\qquad
 1\le r\le r_0,\quad \Lambda=L\Z^2,
\]
where $\|L\|,\|L^{-1}\|\le L_0$ for some unspecified presentation,
$\covol\Lambda\ge V_0$, and different physical bodies have distance
at least $d_0$. The diameter is at most $D_0$, curvature lies in
$[\kappa_-,\kappa_+]$, and each Steiner-centered support function $p_i$
has $\|p_i\|_{C^{6,\beta}}\le M_6$. The type condition is only
\begin{equation}\label{eq:aperture-type}
                   \|p_i-p_j\|_{L^2(0,2\pi)}\ge\sigma\quad(i\ne j).
\end{equation}
The support functions are compared in the common laboratory directions;
there is no minimization over rotations or reflections in this condition.
No rotation or reflection margin, bridge network, protected return
rectangle, onset bracket, or free-area measurement is prescribed.
\end{definition}

Thus symmetric bodies, including disks, are allowed. Rotated congruent
bodies may be different types provided \eqref{eq:aperture-type} holds.
For one type the condition is empty. Every table in the stronger class
of Definition~\ref{def:prior} belongs to this class with the same or weaker
numerical bounds. The type condition implies that $\Lambda$ is the full
translation group: a translation period must carry any chosen component
to a component with the same centered support; it is therefore one of
that component's declared lattice translates. A compact body has no
nonzero translation symmetry.

Put
\begin{equation}\label{eq:aperture-size}
 B=1+L_0,\qquad \tau=\min(1,d_0/8),\qquad
 R=5B+2D_0+3,\qquad Q_R=[-R,R]^2.
\end{equation}
At every attempt choose a uniform position in $Q_R$ and an independent
uniform direction. A solid start is a failure. From a free start keep
only the first impact before time $\tau$, if one occurs, and otherwise
record failure. Only impact positions in $Q_{5B+D_0+1}$ are retained.
Denote the resulting measure on positions together with failure by
$\mathsf F_{\mathcal O}$. The normalization is by the known square area
$4R^2$, not by the unknown free cell area. For finite data the retained
position has a calibrated enclosure of radius at most $\varepsilon_x$.
No later collision, incoming or outgoing angle, boundary parameter,
normal, travel time beyond the cutoff, or obstacle label is needed.
The common laboratory coordinates and resettable launches remain inputs.

\begin{theorem}[Fixed-aperture rigidity and finite acquisition]
\label{thm:aperture-main}
On Definition~\ref{def:aperture-prior}'s class, the single position law
$\mathsf F_{\mathcal O}$ determines the entire periodic obstacle union,
its full translation lattice, the number of types, and its free cell area.
This is an absolute laboratory-frame determination; forgetting that frame
leaves one common Euclidean isometry, type permutation, representative
changes and an integer change of lattice basis.

There are constants $C,c,\nu_0>0$, depending only on the displayed
numerical prior, with the following finite-data version. For
$0<\nu<\nu_0$ and $0<\delta<1/4$, at most
\begin{equation}\label{eq:aperture-cost}
 C\nu^{-3/4}\log\frac{C}{\nu\delta}
\end{equation}
independent attempts from this same square, with
$\varepsilon_x\le c\nu^{3/2}$, suffice with probability at least
$1-\delta$ to recover the correct number of types and the exact rational
relations of a generating period list. The reconstructed supports and
a period basis have error at most $C\nu$ in the finite-presentation
metric defined below, and the free-area estimate has error at most $C\nu$.
Every failed or discarded launch is charged in \eqref{eq:aperture-cost}.
The method is finite: it uses clustering, hulls, smoothing and finite
rational arithmetic, not an optimization over a class of unknown tables.

Neither the aperture nor the time horizon grows with accuracy. No
auxiliary recentering, return selection, endpoint histogram, free-area
normalization, quantitative asymmetry, or clear-bridge prior is used.
The bound counts launch attempts only; localization precision, arithmetic,
travel and apparatus setup are separate resources. No minimax or passive
trajectory assertion is included.
\end{theorem}

The comparison metric uses one representative support function for each
type in its laboratory placement and one basis matrix for the lattice.
It is the infimum of the maximum $C^2$ support error and matrix error
over one common isometry, type permutations, representative translates,
and integer basis changes. It is not the global Hausdorff distance between
two infinite periodic sets with different period matrices. The proof
produces particular bounded representatives and matched bases realizing
the asserted error; a continuously chosen reduced lattice basis is neither
assumed nor needed.

\subsection{All relevant bodies from one cloud}
A first-hit measure is singular relative to plane area. It must not be
estimated by a two-dimensional smooth density. What matters here is its
mass on each short boundary arc.

For any body meeting $Q_{5B}$, its entire boundary lies in
$Q_{5B+D_0}$, and its exterior collar of width $\tau/2$ lies in $Q_R$.
For a boundary arc $I$, parameterize launches by
\[
 q=x-rv,\qquad \tau/4\le r\le\tau/2,\qquad
                 -\pi/3\le\phi\le\pi/3,
\]
where $v$ makes angle $\phi$ with the inward normal at $x\in I$.
The supporting half-plane excludes an earlier hit on that body, and
$r<d_0$ excludes another body. The Jacobian of the first-hit coordinates
is $\cos\phi$. Consequently, with $S=4R^2$,
\begin{equation}\label{eq:aperture-flux}
       \mathsf F_{\mathcal O}(I)\ge\frac{\tau|I|}{8\pi S}.
\end{equation}
This is an unconditional lower bound per attempted launch. It does not
assume a uniform distribution on the unknown boundary.

There are at most
\[
 M=\left\lceil64\left(1+\frac{R+D_0+d_0}{d_0}\right)^2\right\rceil
\]
bodies meeting the relevant square, by packing disjoint $d_0/2$ disks
about points in different bodies. Their perimeters are at most $\pi D_0$
and at least $2\pi/\kappa_+$. Choose
\[
 0<q<\min(d_0/40,1,\pi/\kappa_+),\qquad
 \varepsilon_x<\min(d_0/100,q/10,1/16).
\]
Partition each boundary into arcs of lengths between $q/2$ and $q$.
From \eqref{eq:aperture-flux} a coupon bound gives simultaneous coverage
with failure probability at most $\delta$ as soon as
\begin{equation}\label{eq:aperture-coupon}
 N\ge\frac{16\pi S}{\tau q}
       \log\frac{2M(1+\pi D_0/q)}{\delta}.
\end{equation}
Connect reported points at distance less than $d_0/3$. On this event,
no component joins different bodies, and each relevant body has a
connected complete cloud. Consecutive sampled boundary points have
arclength gaps at most $2q$, whereas different bodies are separated by
$d_0-2\varepsilon_x$. If $P$ is a complete cloud hull, then
\begin{equation}\label{eq:aperture-hull}
 d_H(P,C)=\|p_P-p_C\|_\infty
                 \le e:=\kappa_+q^2/2+\varepsilon_x.
\end{equation}
At a support point the first tangential derivative of the support height
vanishes and the second derivative in arclength is bounded by
$\kappa_+$, which proves this inequality. Localization adds at most
$\varepsilon_x$ to every support value. This is the classical boundary
hull mechanism, here supplied by a physical first-hit minorization;
compare \cite{Schneider,PSSW}.

Only complete clouds may be used for shape recognition. The following
selection avoids assuming that completeness is directly observable.
Write $s(P)$ for the Steiner point of any cluster hull, and keep the
clusters satisfying $\|s(P)\|_\infty\le4B$. Each hull lies in the
$\varepsilon_x$-neighborhood of its true convex body, even for a partial
outer cloud. Since $s(P)\in P$, a kept cluster's body meets
$Q_{4B+\varepsilon_x}\subset Q_{5B}$. The simultaneous event above then
ensures complete coverage and connectivity of that very body. Thus
partial outer components cannot masquerade as additional inner types.
This argument uses raw hulls, not signed smooth approximations, in the
selection test.

\subsection{Recovering the full lattice before any area calculation}
The bounded-presentation hypothesis provides a coverage theorem, not
numerical period labels. For every type, some translate of its Steiner
point lies in $L[-1/2,1/2]^2$, hence has Euclidean norm at most
$L_0/\sqrt2<B$. Since $|s(P)-s(C)|\le2e$, the selected clouds contain
at least one copy of every type when $2e<1/8$.

Temporarily suppose these selected bodies and their centered supports
are known with error at most $u$ in $C^2$, and their centers with error
at most $t$. When $2\sqrt{2\pi}u<\sigma/4$, the threshold $\sigma/2$
on centered $L^2$ support differences separates their exact translation
types: copies agree in the same laboratory directions, while different
types satisfy \eqref{eq:aperture-type}. No orthogonal alignment is
performed. In particular a disk does not cause an orientation singularity.

Choose any cluster with raw center in $Q_{2B}$ as a root. Such a cluster
exists by the preceding coverage argument. Take differences of centers
between all selected copies of its type and the root. Call the true
vectors $d_j$ and their approximations $\widehat d_j$. Every $d_j$ is
a true period. More importantly, if $\ell_1,\ell_2$ are the columns of
the bounded presentation, the root's copies translated by $\ell_i$ have
true centers in $Q_{3B+1/8}$ and estimated centers in $Q_{3B+1/4}$.
They are therefore selected within $Q_{4B}$. Their differences are
$\ell_1,\ell_2$. It follows that
\begin{equation}\label{eq:aperture-generation}
                         \sum_j\Z d_j=\Lambda.
\end{equation}
Any cutoff ambiguity concerns optional additional copies only; it cannot
remove these protected generators or introduce a nonperiod. The bounded
aperture thus certifies saturation and absence of unseen types directly.
There is no area-defect test in this route.

All these differences have norm at most $U=12B$ for small errors.
Set
\begin{equation}\label{eq:aperture-Q}
                     Q=\max\{1,\lceil U^2/V_0\rceil\}.
\end{equation}
Every nonzero determinant of two true periods has magnitude at least
$V_0$, and its perturbation is at most $(2U+t')t'$ when vector errors
are at most $t'$. A threshold $V_0/2$ therefore detects independence
for sufficiently small $t'$. Select a pair with maximal estimated
absolute determinant and write $D=[d_1\ d_2]$. Its index in $\Lambda$
is an integer
\[
                    n=|\det D|/\covol\Lambda\le Q.
\]
Each coordinate of $D^{-1}d_j$ is a bounded rational whose reduced
denominator divides $n$. Distinct rationals with denominator at most
$Q$ differ by at least $Q^{-2}$. Inversion of this $2\times2$ matrix,
using $|\det D|\ge V_0$ and the bound $U$, gives a computable $C_*$
such that coordinate errors are at most $C_*t'$. Once this is less
than $1/(3Q^2)$, finite rational search recovers every coordinate
exactly. One can restrict the numerator to the known coordinate bound
$2U^2/V_0+1$ times $Q$. No integer span of noisy real vectors is taken.

Let $q_*$ be the least common denominator of the recovered coordinates;
it divides $n$ and is at most $Q$. A column Hermite basis $H$ of the
integer group generated by $q_*I_2$ and the vectors $q_*D^{-1}d_j$
then gives
\begin{equation}\label{eq:aperture-lattice}
          \Lambda=(D H/q_*)\Z^2,\qquad
          \widehat L=\widehat D H/q_*.
\end{equation}
The group contains $q_*\Z^2$, so its Hermite determinant is at most
$q_*^2$. All entries of $H$ are bounded by $Q^2$ in the usual
nonnegative remainder convention. This is a finite family, and the
period-basis error is consequently at most $C t'$. Optional repetitions,
different choices of root or independent pair, and cutoff changes only
choose another basis of the same exact lattice. The metric permits
that basis change rather than asserting that a reduced basis varies
continuously.

Keep one body per recovered type. These bodies and
\eqref{eq:aperture-lattice} determine the entire table. The cell's free
area can now be estimated geometrically by
\begin{equation}\label{eq:aperture-area}
              \widehat A=|\det\widehat L|-\sum_i\area(P_i),
\end{equation}
where $P_i$ are the selected raw hulls. On bounded convex bodies the
area error is $O(e)$ under a Hausdorff perturbation $e$, by the parallel
body formula and the uniform perimeter bound. Since the center errors
are $O(e)$, the matrix and free-area errors here are in fact $O(e)$
after type and rational recovery. The $C\nu$ conclusion of the theorem
is a common bound for these errors and the smoother support error.
The scalar estimate \eqref{eq:aperture-area} is not defined by the area
of a separately smoothed numerical curve.

\subsection{Two-derivative confidence and proof of the theorem}
For clarity, we state the smoothing step independently of the return-law
construction. Use the even nonnegative kernel
\[
 \varphi(z)=\frac{315}{256}(1-z^2)^4\one_{[-1,1]}(z),\qquad
 K_h=\tfrac43\varphi_h-\tfrac13\varphi_{2h}.
\]
Its mass is one and its moments of orders one, two and three vanish.
Write $C_\varphi=1+\|\varphi'\|_1+\|\varphi''\|_1\le512$.
For a periodic support function with six derivatives bounded by $B_6$,
\begin{equation}\label{eq:aperture-smoothing}
       \|K_h*p_P-p_C\|_{C^2}
                    \le2C_\varphi e h^{-2}+B_6h^4.
\end{equation}
Differentiating the kernel gives the first term. Taylor expansion of
$p_C^{(j)}$ through degree three, for $j\le2$, gives the second,
since the absolute fourth moment is at most $(20/3)h^4$ and the
Taylor denominator is $24$. These calculations do not differentiate
the noisy polygon. Although $K_h$ is signed, an error less than
$1/(4\kappa_+)$ implies
\[
       K_h*p_P+(K_h*p_P)''>1/(2\kappa_+),
\]
so the reconstructed support still defines a regular strictly convex
body. Convexity is justified by this inequality, not by positivity of
the convolution weights.

Translate by the raw Steiner point before applying this estimate, and
restore the same point afterwards. A uniform $B_6$ follows from $M_6$
and the bounded centering error. The linear Steiner-point functional is
bounded in the support sup norm, so recentering for type comparisons
changes errors by at most a uniform factor. Set, for small $\nu$,
\[
 h=(\nu/(8B_6))^{1/4},\quad
 E_\nu=\frac{\nu h^2}{128C_\varphi},\quad
 q=\sqrt{E_\nu/(4\kappa_+)},\quad
 \varepsilon_x\le E_\nu/8.
\]
Decrease $\nu_0$ to enforce all coverage, selection, convexity, type and
rational thresholds already stated. Then $e\le E_\nu/4$ and
\eqref{eq:aperture-smoothing} is less than $\nu/4$. Deterministic
quadrature can be performed to a further fixed fraction of $\nu$:
the hull has finitely many known linear pieces in its support, the
kernel is explicit, and the integrals have computable bounds. The raw
hull's area and Steiner point can likewise be bounded to a smaller
fraction of $e$. This gives a finite numerical construction with all
arithmetic errors budgeted.

Here $q$ is a constant times $\nu^{3/4}$ and $E_\nu$ a constant times
$\nu^{3/2}$. Substitution into \eqref{eq:aperture-coupon} proves
\eqref{eq:aperture-cost}. On its one simultaneous coverage event, all
subsequent classification, generating-list and rational tests are
correct. No additional probability is spent on data-dependent root or
basis choices: the deterministic argument holds for every choice from
the same covered cloud. The reconstructed finite presentation has
positive curvature and separation for small $\nu$. Only bounded integer
translates can approach a bounded fundamental region, since the matched
lattice basis and its inverse are bounded by the prior and the finite
Hermite family; separation checking therefore has no missing distant
translate. This proves the finite-data assertion.

For exact laws, \eqref{eq:aperture-flux} shows that the support in the
inner region is exactly the union of the complete relevant boundaries.
There are no other impacts, and separation identifies its connected
boundary components and their bounded interiors. Perform the preceding
selection, type identification and period generation without errors.
Every type and a lattice basis are present by the coverage argument.
Equations \eqref{eq:aperture-generation} and
\eqref{eq:aperture-area} then give the exact periodic union and its free
area. This proves the rigidity assertion of
Theorem~\ref{thm:aperture-main}.

\begin{corollary}[Simultaneous certificates without moving the aperture]
\label{cor:aperture-anytime}
Let $\nu_j=2^{-j}\nu_0/2$ and spend
$\delta_j=\delta/[2(j+1)^2]$ on fresh batches from the same square,
using the preceding localization and launch bounds. With probability
at least $1-\delta$ every issued finite-presentation certificate is
correct. By the time target $\nu_j$ is reached, the cumulative number
of attempts is at most $C\nu_j^{-3/4}\log(C/(\nu_j\delta))$.
\end{corollary}
\begin{proof}
The sum of the failure allowances is less than $\delta$ for $j\ge1$.
The simultaneous assertion follows by a union bound. The budgets grow
geometrically, and their logarithmic factors grow at most linearly in
$j$ plus $\log(1/\delta)$; summing gives the displayed bound. Each
batch uses fresh positions with its stated localization, rather than
pretending to refine a previously recorded coarse coordinate. A finite
certificate may be withheld when numerical tests are indeterminate;
on the coverage event the declared thresholds ensure they pass.
\end{proof}

\subsection{Scope and the retained intrinsic law}
For a single disk repeated on a lattice, every rotation and reflection
is a body symmetry, yet the theorem recovers its center translates and
full lattice. Several disks of distinct radii also satisfy
\eqref{eq:aperture-type} whenever the numerical bounds permit their
separated placement. These examples are outside the quantitative
asymmetry assumption of the record-frame theorem. They demonstrate a
strict extension of its geometric scope, not a reformulation of an
asymmetric contact descriptor.

The bounded-presentation hypothesis and translation-type gap are still
substantive. An arbitrarily small local field of view cannot enumerate
unknown periodic types. Nor can arbitrary identical motifs in a
nonprimitive declared cell be assigned an observable cell multiplicity
without further information. We recover the full translation group of
the physical union in the stated type-separated class, not an externally
chosen multiple cell. The position law also retains a common laboratory
frame; it does not solve passive, count-only or marked-length rigidity.

The experiment's reduction is precise: first impacts in one fixed square
are a restriction of the already allowed resettable collision apparatus.
No hypothetical boundary, distance or normalized phase-measure sensor
has been added. The new proof replaces return-law completion by certified
coverage using the existing numerical lattice bound. The intrinsic
identity, local inverse, bridge fingerprint, area-defect certificate,
sequential revalidation and both cloud-cost theorems remain active below.
Their use of the area normalization is indispensable for their
\emph{record-local} information set, not for a globally covering aperture.
The established convex-hull and smoothing rates are ingredients, not a
claim of a new general random-polytope theorem; the conclusion here is
the finite first-impact determination of the complete periodic geometry.

\bigskip\noindent
The remaining sections develop the complementary \emph{record-local}
route. Its stronger prior, auxiliary clouds, two normalized endpoint
histograms and sequential certificates are kept with their full proofs.
Statements about the necessity of area normalization in that route refer
to its local record list, not to the globally covering aperture just proved.
\section{A collision experiment on a finite-smoothness class}
\label{sec:setting}
The local identity behind this paper is a physical density formula, not
a boundary-value measurement. For a selected three-impact return,
\[
 f_T(s,t)=\frac{-W_{st}(s,t)}{2\pi A}(T-W(s,t))_+.
\]
Two active densities determine the stationary action $W$ and the free
cell area $A$. Earlier local recognition used certified boundary and
clearance queries to attach the same intrinsic record to different
physical copies. We now construct the needed geometry from additional
short launches. This changes the acquisition cost and the spatial extent
of the local experiment, while keeping the intrinsic law and its absolute
normalization. In particular, a complete boundary is sampled rather than
continued from a short analytic arc.

\begin{definition}[Numerical prior and directed bounds]\label{def:prior}
Fix $0<\beta\le1$. A table consists of
\[
 \mathcal O=\bigcup_{i=1}^r\bigcup_{\lambda\in\Lambda}(C_i+\lambda),
 \qquad A=\covol\Lambda-\sum_{i=1}^r\area(C_i).
\]
The prior supplies numerical upper bounds
\[
 r_0,D_0,A_1,V_1,b,L_0,M_6,\kappa_+,
\]
and positive lower bounds
\[
 A_0,V_0,a_0,d_0,\kappa_-,\eta,\sigma.
\]
In particular,
\[
 r\le r_0,\quad \operatorname{diam}C_i\le D_0,\quad
 A_0\le A\le A_1,\quad V_0\le\covol\Lambda\le V_1,\quad
 \area(C_i)\ge a_0.
\]
Different physical bodies have distance at least $d_0$. Their curvature
lies in $[\kappa_-,\kappa_+]$. Some presentation $\Lambda=L\Z^2$ has
$\|L\|,\|L^{-1}\|\le L_0$ and a fundamental parallelogram of diameter
at most $b$. No such presentation is supplied. Centered support functions
$p_i$ satisfy $\|p_i\|_{C^{6,\beta}}\le M_6$ and
\begin{align}
 \|p_i(\cdot+\phi)-p_i\|_2&\ge\eta\dist(\phi,2\pi\Z),\label{eq:rot}\\
 \|p_i(-\cdot+\phi)-p_i\|_2&\ge\eta,\label{eq:reflection}\\
 \inf_{O\in O(2)}\|p_i-p_{OC_j}\|_2&\ge\sigma\quad(i\ne j).\label{eq:types}
\end{align}
Choose $2b<R_-<R_+$. Every clear shortest bridge of length below $R_-$
has fixed positive clearance, contact-chart, adjacent-flight, twist,
non-grazing and time-collar margins. Additional admitted bridges have
length at most $R_+$ and the same specified smaller margins. Numerical
lower bounds for these margins are denoted $c_0,r_c,\ell_0,\omega_0,v_*,d_*$,
respectively; numerical bounds for the required chart derivatives are
included. There are fixed nested protected endpoint rectangles on which
$W-2g\le d_*/12$. Active validation times have fixed positive separation.
All bounds are directed: an upper bound may overestimate the true value;
a lower bound may underestimate it, but must be strictly positive.
\end{definition}

One may work in a closed subclass satisfying these bounds with slack
where a protected neighborhood is used. Compactness needs no analytic
strip. Keep bounded presentations $L$ and representative centers $Lt_i$,
$t_i\in[0,1]^2$, instead of choosing a discontinuous reduced basis.
The $C^{6,\beta}$ bound is compact in $C^6$; finitely many body counts
and bounded integer bridge labels suffice on bounded neighborhoods.
Changes at representative faces merely retain another presentation of
the same table. Thus a closed class is compact in the lower norms used
below. The finite-smoothness hypothesis implies uniform $C^{2,\beta}$
endpoint densities: arclength boundary charts are $C^{5,\beta}$,
stationary reduction preserves that regularity for $W$ by the envelope
identity, and the density differentiates $W$ only twice.

The margins in \eqref{eq:rot}--\eqref{eq:types} are still substantive.
They imply that $\Lambda$ is the full translation group: a period must
carry a body to a congruent type, hence to its own declared translate,
and a compact body has no nonzero translational symmetry. This paper
does not infer the numerical priors from the observations. Unlike the
preceding calibration, it requires neither complex-strip bounds nor
local holomorphic graph radii. Smooth nonanalytic perturbations are allowed.

\subsection{The launch and localization contract}\label{sec:sensor}
The apparatus can choose a position uniformly in a specified laboratory
square and choose an independent uniform direction at speed one. It
reports a solid initial position as a failure. At a free position it
records the ordered collisions up to a prescribed bounded time, with
laboratory position enclosures of radius at most a requested tolerance
$\varepsilon_x$. Short auxiliary launches may use translated squares
around the observed contacts and candidate segment. These squares have
a prior-dependent bounded size, large enough to cover an encountered
body of diameter $D_0$, not merely an infinitesimal contact patch.
All launches, including solid positions and trajectories with no collision,
are charged. Position-localization work, travel and setup are separate
from the number of launches. The collision-time and launch controls are
calibrated; no boundary height, derivative, normal, arclength, body label,
period, or distance-to-solid query is supplied.

This is a resettable spatial experiment. Calling it a passive unmarked
trajectory would be incorrect. Its geometric outputs below are confidence
enclosures, with explicitly spent error probability, not deterministic
oracles. The final two windows are two growing two-dimensional endpoint
histograms. They are not two scalar observations.

\begin{theorem}[Collision-generated local geometry]\label{thm:probe-main}
Under the diameter, separation, curvature and $C^3$ consequences of
Definition~\ref{def:prior}, there are computable $C,c,\nu_0>0$ such that
at an encountered bounded flight, $0<\nu<\nu_0$ and $0<\zeta<1/4$,
at most
\begin{equation}\label{eq:probe-main}
 C\nu^{-3/2}\log\frac{C}{\nu\zeta}
\end{equation}
independent auxiliary launch attempts, localized to
$\varepsilon_x\le c\nu^3$, give, with probability at least $1-\zeta$:
complete $C^2$ support-function approximations of the encountered bodies,
local intrinsic coordinates to error $C\nu$, and distance-to-solid
enclosures on a prescribed bounded flight neighborhood. A protected clear
normal bridge is recognized by finite tests with slack. No unknown body
catalogue or analytic continuation is an input to this procedure.
The probability guarantee is uniform conditional on the preceding
history, including the location of the flight.
\end{theorem}

\begin{theorem}[Geometric-size fingerprints]\label{thm:fingerprint-main}
Let
\begin{gather}
 K_0=R_++D_0+1,\qquad C_0=3+9(R_++2D_0)/\eta,\nonumber\\
 \Delta_0=\min\{1,\sigma/36,\eta/36,d_0/(4C_0)\},\quad
 \chi=\Delta_0/4,\quad m=2\lceil32K_0/\Delta_0\rceil.\label{eq:fingerprint-size}
\end{gather}
The gap and the support values of both complete bodies, in their common
normal-contact frame at $m$ equally spaced angles, separate all oriented
bounded bridge classes coexisting in one table by more than $\chi$.
Simultaneous transverse reversal permutes these coordinates.
Theorem~\ref{thm:probe-main} acquires this descriptor to any prescribed
smaller error. Its dimension $1+2m$ is polynomial in the numerical bounds
and reciprocal margins appearing in \eqref{eq:fingerprint-size}; it has
no analytic-continuation exponent. This is not a polynomial bit-complexity
claim and does not identify arbitrary bodies from different tables.
\end{theorem}

\begin{theorem}[Smooth whole-table certification]\label{thm:main}
On Definition~\ref{def:prior}'s class, using only the launch and localization
contract above, a sequential procedure has the following properties.
With probability at least $1-\delta$, $0<\delta<1/4$, every certificate
issued at any epoch is valid for the entire table and its full period
group, up to one common Euclidean isometry, permutation, representative
changes and an integer basis change. Neither a witness graph nor the
number and identities of the types are supplied.

For a chosen $a\ge6$ and target accuracy $\xi>0$, let $T_\xi$ be the
first passing epoch. There are computable class constants $p_*>0,w_0<\infty$
and a deterministic $k_\xi$ such that
\begin{equation}\label{eq:tail-main}
 \Prob(T_\xi>k)\le \frac{\delta}{4(k+1)^a}+w_0e^{-p_*k}
 \quad(k\ge k_\xi).
\end{equation}
In particular $T_\xi<\infty$ almost surely. Writing
\[
 r_k=\frac{\log(C(k+1)^{a+6}/\delta)}{(k+1)^2},
\]
the certified error is at most $C r_k^{\beta/(2\beta+6)}$ for all
sufficiently large passing epochs. No continuation power is applied.
Through epoch $m$ the total number of attempts, including auxiliary
geometric launches, is at most
\begin{equation}\label{eq:cost-main}
 C(m+1)^{11/2}\log\frac{C(m+1)}{\delta}
\end{equation}
with the separately stated rejection-tail probability. Its expectation
up to $T_\xi$ is finite, since $a>11/2$. The largest square half-width is
$O((m+1)^{(\beta+4)/(\beta+3)})$. Localization precision, geometric
postprocessing, travel and setup are not hidden in the launch count.
\end{theorem}

The proof first constructs nearby physical geometry from collision clouds,
then uses the original intrinsic density normalization to exclude missing
types or an unsaturated period subgroup. A finite-smoothness boundary
cannot be inferred from one of its short arcs; no such inference is made.
The full body is sampled in a bounded region. The earlier analytic
short-patch theorem remains a different valid route in the retained volume.

\section{Uniform collision clouds and confidence geometry}\label{sec:clouds}
Put $\tau=\min(1,d_0/8)$, and write $Q_H=z+[-H,H]^2$ for a working
square. For a capped distance radius $c>0$, launch in
$Q_{H+c+2D_0+2\tau}$ and keep first impacts up to time $\tau$.
The square is translated with the observed flight; its area, denoted $S$,
is known. Only first impacts in $Q_{H+c+D_0+\tau}$ are retained.
Every body meeting $Q_{H+c}$ and its outward collar of width $\tau/2$
are contained in the launch square, and its entire boundary is retained.

\begin{lemma}[First-hit minorization]\label{lem:flux-minor}
For every boundary arc $I$ of such a body and every launch attempt,
\begin{equation}\label{eq:minor}
 \Prob\{\hbox{the first impact before }\tau\hbox{ lies in }I\}
                         \ge \frac{\tau|I|}{8\pi S}.
\end{equation}
The bound does not condition on a free initial position or on a collision.
\end{lemma}
\begin{proof}
At $x\in I$ let an incoming direction make angle $\phi\in[-\pi/3,\pi/3]$
with the inward normal, and put $q=x-rv$, $r\in[\tau/4,\tau/2]$.
The segment from $q$ to $x$ lies outside the source body by its supporting
half-plane. It cannot meet another body since its length is below $d_0$.
Thus $x$ is its first impact. The coordinates $(s,\phi,r)$ give phase
volume $\cos\phi\,\dd s\dd\phi\dd r$; the first-hit parametrization is
injective for these states. Dividing by $2\pi S$ gives
$\sqrt3\tau|I|/(8\pi S)$, which implies \eqref{eq:minor}.
Solid launch failures are included in the normalization $S$.
\end{proof}

Only finitely many bodies meet this working region. For example, put
$H'=H+c+3D_0+2\tau+d_0$ and
\[
 M_0=\left\lceil64(1+H'/d_0)^2\right\rceil.
\]
Choose one point from each relevant body; their mutual distances are at
least $d_0$. Disjoint disks of radius $d_0/2$ around these points fit in
the square enlarged by $D_0+d_0$, proving the conservative bound $M_0$.
The circumference of a convex body of diameter $D_0$ is at most $\pi D_0$.

\begin{lemma}[Coverage, clustering and Hausdorff enclosures]\label{lem:cloud}
Choose
\[
 0<q<\min\{d_0/40,1,\pi/\kappa_+\},\qquad
 \varepsilon_x<\min\{d_0/100,q/10\}.
\]
If
\begin{equation}\label{eq:sample-budget}
 n\ge\frac{16\pi S}{\tau q}
       \log\frac{2M_0(1+\pi D_0/q)}{\zeta},
\end{equation}
then, with probability at least $1-\zeta$, the following hold simultaneously.
Every relevant boundary has a retained impact within arclength distance
$q$ of each point. Connect reported positions at Euclidean distance below
$d_0/3$. No component joins different physical bodies, and the cloud of
each fully observed body is connected. Its convex hull $P$ satisfies
\begin{equation}\label{eq:hull}
                      d_H(P,C)\le e_q:=\kappa_+q^2/2+\varepsilon_x.
\end{equation}
Seeds at the encountered collisions may be included with the same
localization tolerance.
\end{lemma}
\begin{proof}
Partition each relevant boundary into arcs with lengths between $q/2$
and $q$; this is possible since its circumference is at least
$2\pi/\kappa_+>2q$. There are at most $1+\pi D_0/q$ arcs per body.
By \eqref{eq:minor}, the probability that any one arc is empty after
$n$ independent attempts is at most $\exp(-n\tau q/(16\pi S))$.
A union bound gives the assertion about coverage. Every point is within
arclength $q$ of a sample in its partition arc; consecutive retained
samples have arclength gaps at most $2q$. Localization therefore still
connects consecutive samples at distance below $d_0/3$. Points from
different bodies remain more than $d_0-2\varepsilon_x>d_0/3$ apart.
Partial clouds from other bodies cannot merge into a complete cloud.

At a support point $x(0)$ in direction $n$, use unit-speed arclength and
$|x''|\le\kappa_+$. For $|s|\le q$,
\[
 0\le n\cdot(x(0)-x(s))\le\kappa_+s^2/2,
\]
since $n\cdot x'(0)=0$. The support function of the exact sample hull
therefore lies between that of $C$ and that function minus
$\kappa_+q^2/2$. Localization changes a hull's support function by at
most $\varepsilon_x$. The support characterization of Hausdorff distance
proves \eqref{eq:hull}, uniformly in direction.
\end{proof}

This argument is a finite-sample form of the boundary random-polytope
mechanism; see \cite{Schneider,PSSW}. What is needed here, and not supplied
by a boundary-sampling oracle, is the physical first-hit minorization,
the treatment of unknown neighboring bodies, and the conditional use of
one finite cloud after an adaptively located flight.

\subsection{No unobserved-solid query}\label{sec:distance}
Take the union $P_{\rm all}$ of convex hulls of all retained components,
including partial exterior components. Let
$d_P(x)=\min\{c,\dist(x,P_{\rm all})\}$, with value $c$ for the empty
union. On the event of Lemma~\ref{lem:cloud}, for every $x\in Q_H$,
\begin{equation}\label{eq:distance-cloud}
 \left|d_P(x)-\min\{c,\dist(x,\mathcal O)\}\right|\le e_q.
\end{equation}
Indeed a hull formed from one body's reported points lies in its
$\varepsilon_x$-neighborhood, so it cannot create a spurious solid farther
than that error from the true solid. If the true distance is below $c$,
a nearest body meets $Q_{H+c}$, is completely covered, and has the
Hausdorff enclosure \eqref{eq:hull}. If the distance is at least $c$,
only the preceding one-sided bound is needed after capping.
Thus \eqref{eq:distance-cloud} controls all query positions at once.
There is no additional pointwise union bound and no distance oracle.

\subsection{Two derivatives and intrinsic coordinates}
Translate a complete cloud by one observed seed. Its true support
function has a $C^3$ bound $B_3=2+M_6+2D_0$, enlarged if necessary to
absorb the seed error. Fix the nonnegative kernel
\[
 \varphi(t)=\frac{315}{256}(1-t^2)^4\one_{[-1,1]}(t),\qquad
 \int\varphi=1,
\]
and its periodic rescaling $\varphi_h$. Set
$\widehat p=p_P*\varphi_h$. This is the support function of an average
of rotations of $P$, and hence remains a convex support function. Let
$C_\varphi=1+\|\varphi'\|_1+\|\varphi''\|_1$; the numerical bound
$C_\varphi\le512$ suffices below. For $h\le1$,
\begin{equation}\label{eq:mollify}
 \|\widehat p-p\|_{C^2}
       \le C_\varphi e_qh^{-2}+B_3h.
\end{equation}
For the first term differentiate the kernel up to twice. For the second,
write $p^{(j)}(\theta-hu)-p^{(j)}(\theta)$ as an integral of
$p^{(j+1)}$, $j\le2$, and use $\int|u|\varphi(u)\dd u\le1$.
This proves the displayed inequality without differentiating a noisy hull.

For $0<\nu<1$ choose
\begin{equation}\label{eq:explicit-cloud-choice}
 \begin{gathered}
 h=\frac{\nu}{8B_3},\qquad
 E_\nu=\frac{\nu^3}{4096C_\varphi B_3^2},\\
 q=\min\{d_0/80,1/2,\pi/(2\kappa_+),
                         \sqrt{E_\nu/(4\kappa_+)}\},\\
 \varepsilon_x\le\min\{E_\nu/8,q/20,d_0/200\}.
 \end{gathered}
\end{equation}
For the small-$\nu$ range where the last entry fixes $q$,
$e_q\le E_\nu/4$. Equations \eqref{eq:mollify} and
\eqref{eq:explicit-cloud-choice} give an error smaller than $\nu/4$.
The finite remaining range is handled by the smaller fixed $q$.
The sample count \eqref{eq:sample-budget} is
$C\nu^{-3/2}\log(C/(\nu\zeta))$, with all constants determined by
bounded local size and the numerical priors. Localization tolerance is
\emph{at most} the displayed value; smaller tolerated errors cost more,
not less, measurement precision.

For small enough $\nu$, $\widehat p+\widehat p''$ stays above
$1/(2\kappa_+)$. The recovered curve
\[
 \widehat x(\theta)=\widehat p(\theta)n(\theta)
                         +\widehat p'(\theta)t(\theta)
\]
is regular and strictly convex. Its radius of curvature differs by
$O(\nu)$ from the true one. Since
$\dd s=(p+p'')\dd\theta$, integration gives intrinsic arclength errors
$O(\nu)$ on bounded arcs. Inverting this positive speed and using the
true $C^3$ bound gives $O(\nu)$ errors for curvature in the common
intrinsic parameter. The Frenet equations then give intrinsic $C^2$
arc error $O(\nu)$. Coordinate rotation to a graph whose tangent has
fixed slack gives $C^2$ graph enclosures with the same order. No
arclength measurement was requested from the apparatus.

All estimates are uniform conditional on the preceding history: after
conditioning, the local square and unknown relevant bodies are fixed,
and only fresh auxiliary launches are used. Quadrature for the convolution
and arc integrals is over explicit hull support functions and a fixed
kernel. Its deterministic errors can be included in the budget
$\nu/4$, since their moduli and the finite hull are known. This proves
the geometric part of Theorem~\ref{thm:probe-main}.

\section{Fourth-order confidence geometry}
\label{sec:compensated}
The positive smoothing kernel in Section~\ref{sec:clouds} uses only a
third-derivative bound. The class of Definition~\ref{def:prior} has six
bounded derivatives. Exploiting those derivatives reduces both the number
of geometric launches and the localization precision required by the
same certificate. No additional observation or regularity assumption is
introduced. The positive-kernel construction remains available under its
weaker local smoothness hypothesis.

Let $\varphi$ be the even kernel in Section~\ref{sec:clouds}, and define
\begin{equation}\label{eq:comp-kernel}
 K_h=\tfrac43\varphi_h-\tfrac13\varphi_{2h},\qquad 0<h\le\tfrac12.
\end{equation}
The periodic rescalings have support in an interval of radius $2h$ before
periodization. This is a signed kernel. Its zeroth moment is one and its
moments of orders one, two and three are zero. Indeed the odd moments
vanish by symmetry, while the second moments of $\varphi_h$ and
$\varphi_{2h}$ are in the ratio $1:4$.

\begin{lemma}[Compensated support approximation]\label{lem:comp-smoothing}
Let $p$ be a periodic support function with $\|p\|_{C^6}\le B_6$, and
suppose $\|p_P-p\|_\infty\le e$. With
$C_\varphi=1+\|\varphi'\|_1+\|\varphi''\|_1$ and
$\widetilde p=K_h*p_P$, one has
\begin{equation}\label{eq:comp-bound}
 \|\widetilde p-p\|_{C^2}
       \le 2C_\varphi e h^{-2}+B_6h^4.
\end{equation}
Here the $C^j$ norm is the maximum of the sup norms of derivatives
through order $j$. If the right side is less than $1/(4\kappa_+)$
and $p+p''\ge1/\kappa_+$, then
$\widetilde p+\widetilde p''>1/(2\kappa_+)$. In particular the
signed correction still defines a regular strictly convex closed curve.
\end{lemma}
\begin{proof}
For $0\le j\le2$,
\[
 \|K_h^{(j)}\|_1\le
 \left(\tfrac43+\tfrac13\,2^{-j}\right)
       h^{-j}\|\varphi^{(j)}\|_1\le 2C_\varphi h^{-2}.
\]
This proves the noise term without differentiating the polygonal support
function. Apply Taylor's formula of degree three to $p^{(j)}$.
The moment identities cancel every nonconstant polynomial term. The
absolute fourth moment is bounded by
\[
 \int |t|^4|K_h(t)|\,\dd t
 \le \left(\tfrac43+\tfrac{16}3\right)h^4
                            \int u^4\varphi(u)\,\dd u
 \le \tfrac{20}3h^4.
\]
The remainder is consequently bounded by
$(20/72)B_6h^4\le B_6h^4$, including $j=2$.
The periodic statement follows by using the periodic lift of $p$ in
this local Taylor calculation. Finally, the errors of $p$ and $p''$
add to at most twice the displayed $C^2$ bound. The positive lower bound
for the radius of curvature proves strict convexity and regularity.
Convexity here follows from this inequality, not from averaging convex
bodies with positive coefficients.
\end{proof}

\begin{theorem}[Reduced geometric launch budget]\label{thm:comp-probe}
Use the full $C^{6,\beta}$ class of Definition~\ref{def:prior}. The
outputs and conditional confidence guarantee of
Theorem~\ref{thm:probe-main} can instead be obtained with
\begin{equation}\label{eq:comp-probe}
 C\nu^{-3/4}\log\frac{C}{\nu\zeta}
\end{equation}
auxiliary launch attempts and impact-position error tolerance at most
$c\nu^{3/2}$. The constants depend only on the same numerical priors.
Neither an analytic continuation constant nor a boundary-query oracle
is used. Deterministic convolution and quadrature errors are included
in the geometric error budget; their arithmetic work is not included
in the launch count.
\end{theorem}
\begin{proof}
Translate the cloud by one of its measured seeds. One may take
$B_6=2+M_6+2D_0$ for the true support in this frame. For all sufficiently
small $\nu$, use
\begin{equation}\label{eq:comp-choices}
 \begin{gathered}
 h=(\nu/(8B_6))^{1/4},\qquad
 E_\nu=\frac{\nu h^2}{128C_\varphi},\\
 q=\min\{d_0/80,1/2,\pi/(2\kappa_+),
                         \sqrt{E_\nu/(4\kappa_+)}\},\\
 \varepsilon_x\le\min\{E_\nu/8,q/20,d_0/200\}.
 \end{gathered}
\end{equation}
The remaining bounded range uses the smaller fixed thresholds. The
confidence hull of Lemma~\ref{lem:cloud} has error
$e_q\le E_\nu/4$. Equation~\eqref{eq:comp-bound} gives
\[
 2C_\varphi e_qh^{-2}+B_6h^4
                 \le\nu/256+\nu/8.
\]
Allow an additional $\nu/16$ for numerical convolution and integration.
The total is below $\nu/4$. For small $\nu$ the preceding lemma
therefore supplies a strictly convex support function, and the intrinsic
coordinate, normal-root, fingerprint and clearance arguments apply with
the same errors. Distance-to-solid enclosures continue to use the original
hulls and \eqref{eq:distance-cloud}, rather than an unjustified
monotonicity of the signed reconstruction.

In the asymptotic range $E_\nu=c_1\nu^{3/2}$ and
$q=c_2\nu^{3/4}$. Substitution into the exact first-hit coverage budget
\eqref{eq:sample-budget} proves \eqref{eq:comp-probe}. These are
attempts, including solid initial positions and no-impact trajectories.
The confidence argument is unchanged after conditioning on any preceding
history: the square and bodies are fixed and the auxiliary launches are
fresh. Every new smoothness bound used here is already supplied by
Definition~\ref{def:prior}. No differentiability of the noisy hull is
assumed.
\end{proof}

\begin{corollary}[Reduced total launch cost]\label{cor:comp-work}
In Theorem~\ref{thm:main}, replace each positive-kernel cloud call by
Theorem~\ref{thm:comp-probe}, keeping the same raw proposals, fresh
validation, error spending and schedule \eqref{eq:schedule}.
The anytime certificate, the stopping tail \eqref{eq:tail-main} and
the certified reconstruction error are unchanged. The cumulative launch
bound improves from \eqref{eq:cost-main} to
\begin{equation}\label{eq:comp-work}
 C(m+1)^{19/4}\log\frac{C(m+1)}{\delta},
\end{equation}
with the same separate rejection-tail allowance. Impact localization
need only satisfy $\varepsilon_{x,k}\le c\nu_k^{3/2}$.
The expected total number of launches up to $T_\xi$ is finite whenever
the spending exponent satisfies $a>19/4$ and the retained validity
conditions hold; in particular every stipulated $a\ge6$ suffices.
\end{corollary}
\begin{proof}
There are at most $C_1(k+1)^3$ clouds at epoch $k$.
The schedule gives $\nu_k^{-3/4}\le C(k+1)^{3/4}$ in the small-error
range, and
$\log(C/(\nu_k\zeta_{k,j}))\le C\log(C(k+1)/\delta)$.
Thus the epoch budget is
$C(k+1)^{15/4}\log(C(k+1)/\delta)$, including the smaller expected
primary-rejection cost. Summation gives \eqref{eq:comp-work} and the
same rejection event as before. No probability argument in
Section~\ref{sec:sequential} depends on the internal smoothing choice:
it uses only genuine uniform geometric enclosures, conditional failure
allowances and deterministically capped call counts, all retained here.

The event that epoch $k$ is executed is measurable before its fresh
randomness. Multiplying the conditional expected epoch cost by the tail
bound at $k-1$ gives a convergent series because
$15/4-a<-1$. The exponentially decaying witness term is summable as well.
The finitely many early epochs contribute a finite amount. Localization
work, apparatus travel and geometric postprocessing remain distinct from
the launch count; no polynomial bit-complexity or minimax claim follows.
\end{proof}

The order reduction is a consequence of the already stated smoothness,
not a stronger identifiability assumption. In particular the original
third-derivative confidence construction and all of its bounds remain
valid. The sharper implementation records explicitly where the six
bounded derivatives improve acquisition rather than being used only for
the regularity of the endpoint density.

\section{Finite recognition and protected geometric tests}\label{sec:recognition}
We distinguish a statistical reconstruction of a physical body from its
translation class in the unknown periodic table. Spatial clustering in
Lemma~\ref{lem:cloud} solves the first problem locally. The following
finite descriptor solves the second without extending a contact germ.

\subsection{Support fingerprints in the normal frame}
For an oriented shortest bridge, translate its source contact to zero
and rotate the normal towards the target to $(1,0)$. The transverse axis
has either of two signs. Write $p_0^e,p_1^e$ for the support functions of
the two complete bodies in this one frame, without separate recentering.
Their absolute values and first derivatives are bounded by $K_0$ in
\eqref{eq:fingerprint-size}. Define
\begin{equation}\label{eq:descriptor}
 F(e)=\bigl(g_e,(p_0^e(2\pi j/m))_{j=0}^{m-1},
                    (p_1^e(2\pi j/m))_{j=0}^{m-1}\bigr).
\end{equation}
Reflection across the normal permutes $j$ to $-j$ modulo $m$ in both
blocks. It is a coherent reversal of the pair, not separate endpoint signs.

\begin{lemma}[Locking physical copies from complete pairs]\label{lem:locking}
Suppose two normal-frame pairs from one table have support differences
at most $\Delta_0$ on the full circle. Then they represent the same
oriented bridge up to translation.
\end{lemma}
\begin{proof}
If two support functions differ by at most $\Delta$, their Steiner
points differ by at most $2\Delta$, since
$s(C)=\pi^{-1}\int_0^{2\pi}p_C(\theta)n(\theta)\dd\theta$.
Their centered support functions consequently differ by at most
$3\Delta$ and by at most $9\Delta$ in $L^2$. The shape gap excludes
different source types when $\Delta\le\Delta_0$. The exact congruence
between the two source copies, expressed in their normal frames, has
orthogonal part $O$. The reflection margin excludes its orientation-reversing
component; the rotation margin bounds its angle by $9\Delta/\eta$.

Align the two physical source copies by their true lattice translation.
The preceding centering and angle estimates show that their target
images then differ in Hausdorff distance by at most
\[
 \{3+9(R_++2D_0)/\eta\}\Delta=C_0\Delta<d_0.
\]
For completeness, the centering displacement costs $2\Delta$; rotation
about a source Steiner point costs at most
$(R_++2D_0)9\Delta/\eta$ on the pair, and the original target comparison
costs $\Delta$. Distinct target bodies in that physical table have
separation at least $d_0$, so the two targets coincide. The closest
segment of two disjoint strictly convex bodies is unique. It therefore
fixes both contacts, the normal and, by the reflection exclusion, the
transverse sign. This proves the assertion.
\end{proof}

\begin{proof}[Proof of Theorem~\ref{thm:fingerprint-main}]
If the node differences are at most $\chi=\Delta_0/4$, the derivative
bound gives, at any circle point,
\[
 |p_b^e-p_b^{e'}|\le\chi+2K_0\pi/m
                          \le\Delta_0/4+\Delta_0/8<\Delta_0.
\]
Here $\pi<4$ and $m\ge64K_0/\Delta_0$ were used. Apply
Lemma~\ref{lem:locking}. The contrapositive gives separation greater
than $\chi$ for inequivalent oriented classes. All ingredients of
\eqref{eq:fingerprint-size} are finite arithmetic in numerical margins.
With each estimated descriptor within $\chi/16$ of its truth, same-class
distances are at most $\chi/8$, and different-class distances exceed
$7\chi/8$. A threshold $\chi/2$, testing both coherent reversals,
therefore recognizes the class and sign. Fixed-order support, normal
and gap errors from Theorem~\ref{thm:probe-main} make these comparisons
valid once $C\nu<\chi/16$.
\end{proof}

This polynomial-size fingerprint uses whole-body support values produced
by bounded-range sampling. It is not a shorter list of values on the
same short contact interval used in the preceding analytic theorem.
That distinction explains why its separation needs no continuation
constant. It also prevents a false claim that a smooth body is globally
determined by a finite contact jet.

\subsection{One explicit normal-root certificate}\label{sec:root}
We give quantitative choices for the local root and clearance tests.
In true unit-speed charts near a closest pair, $D(s,u)=|x(s)-y(u)|$
satisfies
\begin{equation}\label{eq:normal-hessian}
 \nabla D=0,\qquad D''=
 \begin{pmatrix}g^{-1}+\kappa_0&-g^{-1}\\
                -g^{-1}&g^{-1}+\kappa_1\end{pmatrix}.
\end{equation}
Its least eigenvalue is at least $\mu=\kappa_-$. In unit-speed charts
one may use the explicit Hessian-Lipschitz bound
\[
 L_D=64\{1+d_0^{-2}+\kappa_+/d_0+\kappa_+^2
                                      +2M_6\kappa_+^3\}.
\]
Indeed, writing $\rho=p+p''$, one has
$|\dd\kappa/\dd s|=|\rho'|/\rho^3\le2M_6\kappa_+^3$.
The first three derivatives of a unit-speed arc are bounded by
$1$, $\kappa_+$ and $2M_6\kappa_+^3+\kappa_+^2$.
The first three derivatives of Euclidean distance at separation at least
$d_0$ are bounded by $1,d_0^{-1},6d_0^{-2}$. The multivariate chain
rule gives the conservative displayed bound.
Set
\begin{equation}\label{eq:root-radius}
 r=\min\{r_c/8,d_0/32,\mu/(32L_D)\}.
\end{equation}
The initial encountered patches are accepted only inside a protected
neighborhood on which the true Hessian has least eigenvalue at least
$3\mu/4$. This condition is verifiable with the cloud's derivative
enclosures; genuine witness returns have uniform slack.

Here is a specific certificate. At a trial center $z$, let $B$ be a
symmetric approximation to $D''(z)$ with error at most $\mu/16$ and
least eigenvalue at least $\mu/2$. Let $\epsilon_g$ bound the gradient-approximation error uniformly
on the ball, as provided by the cloud enclosures. Verify
the inequality
\begin{equation}\label{eq:newton-test}
 \|\widehat{\nabla D}(z)\|+\epsilon_g\le\mu r/8,
 \qquad \|D''(x)-B\|\le\mu/4\quad(x\in\overline B(z,r)).
\end{equation}
The second condition follows from $L_Dr+\mu/16\le\mu/4$ when the
ball stays in the certified rectangle. The map
$x\mapsto x-B^{-1}\nabla D(x)$ has Lipschitz constant at most $1/2$
on this ball and moves its center by at most $r/4$. It maps the ball
strictly into itself and has a unique fixed point. The root lies within
$2\|B^{-1}\|\epsilon_g$ of the corresponding root for the central
approximation, with the approximation residual added to $\epsilon_g$
when it is not zero. This gives finite root enclosures, rather than an
unstated interval-Newton call.

Approximate roots can be found by a finite mesh followed by contraction,
or by minimizing distance between the reconstructed strictly convex
bodies. The uniform Hessian and chart slack make the tests pass for all
sufficiently accurate clouds of a protected bridge. Refinement is capped
at the deterministic sample and arithmetic budgets; indeterminate inputs
are rejected. Convex supporting half-planes show that a certified facing
normal is the global shortest segment for those two bodies.

\subsection{Clearance and record assignment}
Let $C_0^e,C_1^e$ denote the two \emph{physical bodies} defining the
candidate bridge; different copies of the same type are still different
physical bodies. Define
\[
 \mathcal O_e^{\rm other}
  =\bigcup_{C\text{ a physical body},\ C\ne C_0^e,C_1^e} C.
\]
Delete exactly the two corresponding spatial cloud components from
$P_{\rm all}$, not every component with either endpoint's shape. The
argument for \eqref{eq:distance-cloud}, applied after this deletion,
gives the same capped-distance enclosure for $\mathcal O_e^{\rm other}$.
The padded cloud region must contain every body meeting the capped
neighborhood of the entire tested segment. Partial outer components are
kept, just as in Section~\ref{sec:distance}.

Put $\bar c=\min(c_0,d_0/4)$ and use cap at least $\bar c$. On the
closed candidate segment, a mesh of spacing at most $\bar c/8$, total
distance and root-position enclosure radius at most $\bar c/8$, and
reported distances at least $3\bar c/4$ imply
\[
 \dist([x_0,x_1],\mathcal O_e^{\rm other})\ge\bar c/2.
\]
This follows directly from the one-Lipschitz property of distance.
Protected bridges with the stipulated clearance pass at sufficiently
small error. The endpoint bodies have deliberately not been included
in this inequality.

Their contribution is separate and exact. Write $x_1=x_0+gn$ for the
certified common normal. The facing supporting half-planes give, for
$0\le s\le g$,
\[
 \dist(x_0+sn,C_0^e)=s,\qquad
 \dist(x_0+sn,C_1^e)=g-s.
\]
Choose $a=\min(d_0/8,g/4)$ and treat $s\in[0,a]$ and
$s\in[g-a,g]$ in the protected endpoint charts, restricted to the
exterior sides of the supports. On $[a,g-a]$ a tube of radius smaller
than $\min(a/2,\bar c/4)$ avoids all solids. Root uncertainty is included
in the enclosure budget above. An upper bound on $a$ is not used as a
lower bound on endpoint distance; the explicit tube radius uses $a$
itself. All distances are computed from this same cloud, not obtained
from a separate geometric sensor.

At most one record is extracted from a trace, and it must have exactly
three impacts before its recorded horizon. First and third impacts must belong to the same spatial cloud component,
with the middle impact on the other body. Inside fixed protected patches,
the stationary two-flight branch is unique. Record assignment additionally
checks the collision order, gap cutoff below $R_+$, protected patches,
twist and adjacent-flight margins. Every trajectory in a common smaller
return rectangle of a clear bridge of length below $R_-$ passes these
checks on the confidence event. The slack $R_+-R_-$ protects the cutoff.
For a registered branch, future trajectories are classified with the
same support descriptor and assigned intrinsic coordinates using the
reconstructed charts and their true-error enclosures. Interior histogram
rectangles have a fixed collar. Ambiguity in a bin is consequently only
an endpoint-coordinate band of thickness $C\nu$, not an uncharged change
of branch. Failed or out-of-rectangle trajectories remain outcomes.

This proves the record part of Theorem~\ref{thm:probe-main}. Its confidence
event is simultaneous for all geometry extracted from one cloud. It does
not say that each statistically accepted record is infallible. The next
two sections explain the finite-list certificate and how independent
revalidation makes the distinction harmless for an anytime guarantee.

\section{Intrinsic normalization and smooth completion}\label{sec:certificate}
The full-body cloud changes one step of the earlier inverse: complete
images now have measured confidence bounds. Analytic continuation is not
needed. The boundary laws retain their essential role in calibrating the
free area of the unknown periodic cell, including unseen obstacles.

\subsection{The retained local inverse}
For facing unit-speed charts put
\[
 D(s,u)=|x(s)-y(u)|,\quad
 W(s,t)=D(s,u(s,t))+D(t,u(s,t)),\quad D_u(s,u)+D_u(t,u)=0.
\]
The positive intermediate Hessian selects the physical branch. In
collision-section coordinates momentum is $-W_s$ and Liouville volume is
$(-W_{st})\dd s\dd t\dd r$. On a collar shorter than adjacent flights,
$0<r<T-W$ is precisely the allowed residual interval before the first
impact. Thus
\begin{equation}\label{eq:law}
 f_T=\frac{-W_{st}}{2\pi A}(T-W)_+.
\end{equation}
At two strictly active times $T_2>T_1$,
\begin{equation}\label{eq:action-area}
 W=\frac{T_1f_2-T_2f_1}{f_2-f_1},\qquad
 A=\frac{(T_2-T_1)(-W_{st})}{2\pi(f_2-f_1)}.
\end{equation}
This is the same intrinsic observation as before. It is not conditioned
on return success. The two density functions carry the absolute
normalization; the finite fingerprint does not replace them.

For reference, on the diagonal let
$\ell=W(s,s)/2$, $a=W_s(s,s)$ and $v=\sqrt{1-a^2}$. With source tangent
$T$ and inward normal $N$, the nearest-point direction is $aT+vN$.
Differentiating distance and eliminating the stationary intermediate point
by its Schur complement gives
\begin{align}
 \kappa(s)&=\frac{W_{ss}-W_{st}-v^2/\ell}{v},\label{eq:curvature}\\
 \kappa_o(u(s))&=\ell^{-1}\left(\frac{v^2}{-2\ell W_{st}}-1\right),
 &u'(s)&=\frac{v}{1+\ell\kappa_o(u(s))}>0.\label{eq:foot}
\end{align}
Indeed $D_{ss}=v^2/\ell+\kappa v$, $D_{su}=-v/\ell$ and
$D_{uu}=1/\ell+\kappa_o$, while
$W_{st}=-D_{su}^2/(2D_{uu})$ and $W_{ss}=D_{ss}+W_{st}$.
Frenet integration determines the source arc, and the target is
$y(u(s))=x(s)-\ell(aT+vN)$.

\begin{proposition}[Local physical stability]\label{prop:local-inverse}
Assume uniform physical $C^3$ bounds on larger arcs, positive lower gap,
area, curvature, active, twist, time-separation, adjacent-flight and
clearance bounds, a diagonal bound $v\ge v_*>0$, and uniformly bounded
$C^2$ densities with $f_2-f_1$ bounded below. A $C^2$ perturbation of
size $e$ between two physical density pairs changes the gap, free area
and both intrinsic $C^2$ arcs by at most $Ce$, up to common reflection.
The same area-error bound holds for a numerical density pair within $e$
of a physical one, provided the denominator is kept positive.
\end{proposition}
\begin{proof}
Division is Lipschitz on the stated bounded $C^2$ sets. Apply
\eqref{eq:action-area} and \eqref{eq:curvature}--\eqref{eq:foot}.
For the opposite arc first integrate $u'$ and compare inverse parameters;
the physical $C^3$ bound controls curvature under that parameter change.
Frenet stability gives the $C^2$ arcs. This uses no third action derivative.
The area estimate is simply a smooth rational operation in the value,
first and second derivatives of the recovered action and the densities.
\end{proof}

\subsection{Noisy full images and the period group}
Suppose every entry in a finite genuine record list has complete-pair
support error at most $u$ in $C^2$, in its contact frame, and density
error at most $e$. Cluster the centered images modulo $O(2)$ at threshold
$\sigma/2$. Small enough $u$ gives the exact visible types by
\eqref{eq:types}. A record is an edge of the resulting multigraph;
loops and repeated records are retained.

Place a spanning tree in a component. Approximate matches of the same
source image have the same reflection component and rotation difference
$O(u/\eta)$ by \eqref{eq:rot}--\eqref{eq:reflection}. Steiner points
control translation. The $C^3$ bound controls the action of a small
rotation on $C^2$ supports. Tree paths have at most $r_0-1$ edges, so
all placed representatives and each non-tree cycle displacement $d_e$
have error $C u$, independently of repetitions. Put, conservatively,
\[
 B_0=4r_0(R_++2D_0),\qquad Q=\max(1,\lceil B_0^2/V_0\rceil).
\]
All true cycle displacements have norm at most $B_0$ and lie in $\Lambda$.
Let $I$ be the visible types in the component and
$\Gamma=\sum_e\Z d_e$. A rank-two component has completion defect
\begin{equation}\label{eq:defect}
 \mathcal D=\covol\Gamma-A-\sum_{i\in I}\area(C_i)
 =([\Lambda:\Gamma]-1)\covol\Lambda+
                         \sum_{i\notin I}\area(C_i).
\end{equation}
Consequently $\mathcal D=0$ exactly for a complete saturated component;
otherwise it is at least $g_0=\min(V_0,a_0)>0$.

We recall why noisy cycles can be used without taking integer spans of
noisy real vectors. If their errors are at most $t=C u$, a pair determinant
has error at most $(2B_0+t)t$. A nonzero true determinant is at least
$V_0$. Once this error is below $V_0/4$, threshold $V_0/2$ recovers rank.
Choose a pair with maximal estimated determinant and denote the true pair
by $D=[d_1\ d_2]$. Then
\[
 V_0\le|\det D|\le B_0^2,\qquad
 n_D=|\det D|/\covol\Lambda\in\{1,\ldots,Q\}.
\]
Each coordinate of $D^{-1}d_e$ is a bounded rational with denominator
dividing $n_D$. Distinct reduced fractions with denominators at most $Q$
are at least $Q^{-2}$ apart. Matrix inversion with the preceding bounds
therefore permits a finite rational search that recovers every coordinate
exactly for $t<c/Q^2$. Let $q\le Q$ be their least common denominator.
A column Hermite basis $H$ for the integer group generated by $qI_2$
and all $qD^{-1}d_e$ gives
\begin{equation}\label{eq:hermite}
 \Gamma=(D/q)H\Z^2,\qquad
 \covol\Gamma=|\det D|\,|\det H|/q^2.
\end{equation}
The possible Hermite forms are finite for $q\le Q$, and
$1\le|\det H|\le q^2$. Thus the covolume error is $O(u)$ after
integer recovery. No noisy integer span is formed.

The area functional
$\area(C)=\tfrac12\int_0^{2\pi}(p^2-(p')^2)\dd\theta$
is Lipschitz on bounded $C^1$ sets. Use one area per visible type, the
common free-area estimate in \eqref{eq:action-area}, and
\eqref{eq:hermite}. The defect error is at most $C(u+e)$.

\begin{theorem}[Finite-smoothness completion certificate]\label{thm:certificate}
There are computable $u_*,e_*,C>0$ depending only on the numerical priors
such that, for any finite list of genuine records with simultaneous
full-pair error $u<u_*$ and density error $e<e_*$, visible incidence and
rational cycle coordinates are correctly recovered. If
$C(u+e)<g_0/4$, accept a rank-two component precisely when
$|\widehat{\mathcal D}|<g_0/2$. Every accepted component is complete and
saturated, and every complete saturated component passes. The support
and period-basis error is at most $Cu$, in particular at most $C(u+e)$.
The same event controls every component, repetition and selected deletion.
\end{theorem}
\begin{proof}
The preceding estimates prove all the discrete statements. The positive
defect gap proves the acceptance equivalence. A recovered period basis
is $\widehat D H/q$, compared with $DH/q$; the finite Hermite family
bounds its error by $Cu$. The placed-body errors have that same bound.
For sufficiently small $u$, curvature and physical separation persist;
only bounded integer translates can approach a bounded fundamental
region, so there is no unexamined collision with a distant translate.
All operations concern finitely many approximated supports, compact
one-dimensional rotation searches and finite rational searches. A grid
search with a known derivative bound computes the matching residual to
any required smaller error. No fit over analytic functions or continuation
is invoked. Uniform constants do not count repetitions, since tree depth
and the number of visible types are bounded by $r_0$.
\end{proof}

The normalized boundary law still does indispensable work here. Local
clouds near the records need not enumerate a period cell. Their visible
areas alone cannot certify that unseen types or extra periods are absent.
The free area supplied by the two unconditional density functions is
exactly the missing term in \eqref{eq:defect}. Conversely, claiming smooth
whole-body recovery from the density's short arc alone would be invalid;
the complete-pair confidence input comes from actual first-hit sampling.

\section{Fresh validation and finite expected acquisition cost}\label{sec:sequential}
A statistically wrong descriptor cannot be treated as a permanent name.
This is the main distinction between the new confidence probe and the
previous deterministic local oracle. We retain raw proposals and rebuild
all geometric decisions at every epoch.

\subsection{A witness exists without being supplied}
Take $b<c<R_-/2$. The open $c$-neighborhoods of all physical bodies
cover the plane: a translate of any fixed obstacle point is within $b$
of any point by the fundamental-parallelogram diameter bound. The
intersection graph of this open cover is connected. Its adjacent bodies
have gap below $R_-$. If a shortest segment between such bodies is
blocked, an intervening body has strictly smaller gap to each endpoint
body. There are only finitely many gap values below $R_-$ modulo periods,
because bounded presentations give bounded integer copy labels.
Induction on these gap values replaces each blocked edge by a path of
clear shorter edges. Hence the lifted clear-bridge graph is connected.
Its quotient visits every type, and paths to every lattice translate
show that its cycle displacements generate the full $\Lambda$.

Keep a quotient spanning tree. Two independent cycle vectors have index
at most $Q$ in $\Lambda$, by the bound preceding \eqref{eq:defect}.
Whenever the kept subgroup is proper, another fundamental cycle reduces
that index to a proper divisor, at most half its previous value. Thus
there is a saturated witness with at most
\begin{equation}\label{eq:witness}
                     w_0=r_0+1+\lfloor\log_2 Q\rfloor
\end{equation}
records. Neither its membership nor its combinatorial graph is supplied.
The argument is entirely geometric and uses no analyticity.

\subsection{Laboratory preparation and raw discovery}
For primary launches choose a uniform point in $[-L,L]^2$ and reject
solid positions before choosing an independent direction. Rejections
are counted. For $L\ge2b$, the induced free-position law modulo the
unknown lattice differs in total variation from quotient free measure by
\begin{equation}\label{eq:square}
 \tau_L\le\min\{1,8V_1b/(A_0L)\},\qquad
 \Prob\{q\hbox{ free}\}\ge\pi_0=A_0/(4V_1).
\end{equation}
Indeed the union of full fundamental tiles contained in the square
covers its inner square of half-width $L-b$; the remaining area is at
most $8Lb$. The free part of the full tiles projects exactly to the
quotient law, and their free fraction is at least $A_0/V_1$.
These comparisons prove both inequalities without knowing any tile.
Auxiliary cloud launches instead use \emph{all} attempted positions and
the known local-square area, as in Lemma~\ref{lem:flux-minor}.

Choose a fixed time grid in $[0,2R_-+d_*]$ fine enough that every witness
gap $g$ has a grid time with $T-2g\in[d_*/3,d_*/2]$. On a protected
rectangle $[-a_0',a_0']^2$ small enough that $W-2g\le d_*/12$,
\eqref{eq:law} has a uniform lower bound. Denote the resulting lower
return probability by $\mu_*>0$. For example a smaller positive number
than $\omega_0d_*(a_0')^2/(2\pi A_1)$ suffices. Choose a fixed discovery
square with $\tau_L\le\mu_*/2$. At each epoch retain the whole raw
bounded-time trace from one accepted free launch at a uniformly chosen
grid time. The raw positions need only fixed coarse enclosures small
relative to the protected chart radius; they serve as seeds for later
clouds, not as exact histogram coordinates. Each missing witness has
conditional raw-arrival probability at least
\begin{equation}\label{eq:hazard}
                         p_*=\mu_*/(2H_0)>0,
\end{equation}
where $H_0$ is the grid size. No acceptance test is used to discard these
raw proposals. If $K_{\rm raw}$ is the first epoch whose raw list contains
the witness, then
\begin{equation}\label{eq:coupon}
                    \Prob(K_{\rm raw}>k)\le w_0e^{-p_*k}.
\end{equation}
Repeated entries do not affect this estimate.

\subsection{Fresh geometric and density validation}
At epoch $k$ independently reconstruct clouds for each of the $k$ raw
proposals. Reject those without a protected genuine branch certificate;
form a new registry from the others. No old class assignment, sign,
geometry or rejection decision is reused. Repeated valid records may all
be kept, so the current length $J_k$ is at most $k$. Coarse old seeds
suffice because their fixed neighborhood contains the protected root;
new clouds, not impossible refinements of old measured positions, supply
the new accuracy. Old coarse points are not inserted into a finer convex
hull. With their fixed error below $d_0/20$, they are associated to the
unique cloud component within that error plus the current Hausdorff bound;
the $d_0$ separation excludes a second such body. Raw acceptance uses
interval tests with slack on these coarse enclosures. Protected witness
returns lie strictly inside those tests.

Conditional on these proposal clouds, choose two active times for each
current record. For a rough gap estimate with error below $d_*/100$, one
may use $2\widehat g+d_*/3$ and $2\widehat g+d_*/2$. Draw
$N_k=(k+1)^2$ fresh accepted free trajectories at each time. Any needed
record classification and endpoint arclength are computed from additional
fresh clouds around those trajectories, at the current accuracy $\nu_k$.
They are not borrowed from their density samples. The total number of
cloud calls is at most $C_1(k+1)^3$, since the time horizon and the number
of impacts per trace are bounded by separation. All horizons, candidate
counts and refinement budgets are deterministically capped even when
a confidence event fails. In particular no erroneous classification
can create an unbounded work request.

Let $a\ge6$ and define
\begin{equation}\label{eq:spending}
 \alpha_k=\frac{\delta}{32(k+1)^a},\qquad
 \zeta_{k,j}=\frac{\alpha_k}{2C_1(k+1)^3}
       \quad(1\le j\le C_1(k+1)^3).
\end{equation}
Each cloud is sampled with that conditional error allowance and its
deterministic budget \eqref{eq:sample-budget}. Uniform conditional
coverage and iterated conditioning give an event $G_k$ on which all
these geometric enclosures hold, with $\Prob(G_k^c)\le\alpha_k/2$.
The count is deterministic even when the actual locations and calls are
adaptive. Off this event the procedure may make a wrong provisional
decision; that decision is not carried into a later epoch.

For the concentration step, one must not condition the validation
trajectories on their cloud-success event. Instead compare their measured
cell indicators with the true indicators and the true boundary-band
indicators. With a bin of side $h$ and coordinate error at most $C\nu$,
a changed assignment lies in a band of area at most $C(h\nu+\nu^2)$.
On the geometric event, these indicators dominate every possible
numerical or cloud-based bin change pathwise. For a bounded physical
density the true bin probability is at most $C h^2+\tau_L$ and the
band probability is at most $C h\nu+\tau_L$, provided $\nu\le c h$.
Bernstein's inequality for the true bin and band indicators, conditional
only on the proposal history and chosen times, gives simultaneous cell
error
\begin{equation}\label{eq:cell}
 C\{h\sqrt{r_N}+r_N+\tau_L+\nu h\},\qquad
 r_N=\frac{\log(C(J\vee1)N/\alpha_k)}{N},
\end{equation}
except on an additional event of probability at most $\alpha_k$ on
good proposal histories. This last formulation suffices: a bad proposal
history is already charged to $G_k^c$. A union bound over the
$O(Jh^{-2})$ bins is absorbed in the displayed logarithm when $h^{-2}\le CN$.
One endpoint outcome gives all its bin indicators simultaneously.

A $3\times3$ cell-average stencil reproduces polynomials of coordinate
degree at most two: in one coordinate its matrix has rows
$(1,i,i^2+1/12)$, $i=-1,0,1$, and is invertible. Taylor remainder bounds
for $C^{2,\beta}$ densities, followed by a bounded-overlap smooth
partition of unity, therefore give a $C^2$ density error
\begin{equation}\label{eq:density}
 C\{h^\beta+\sqrt{r_N}h^{-3}+(r_N+\tau_L)h^{-4}+\nu h^{-3}\}.
\end{equation}
For clarity a mass error is divided by cell area $h^2$, and two further
derivatives cost $h^{-2}$. The bias of order $j\le2$ is
$O(h^{2+\beta-j})$; partition derivatives consume the corresponding
lower-order powers. No derivative estimate is taken directly from
total variation.

Use the deterministic larger value $r_k$ in Theorem~\ref{thm:main} and
\begin{equation}\label{eq:schedule}
 h_k=r_k^{1/(2\beta+6)},\quad L_k\ge C h_k^{-(\beta+4)},\quad
 \nu_k=c\min\{\nu_0,\chi,r_k^{1/2}\}.
\end{equation}
Constants include the coordinate conversion and descriptor errors.
Then $h_k^{\beta+3}=r_k^{1/2}$, and every term of \eqref{eq:density}
is at most $C r_k^{\beta/(2\beta+6)}$. Complete-pair cloud errors are
$C\nu_k$, which are smaller in the small-error range. Early epochs
outside that range collect data but issue no certificate.

\subsection{Anytime validity, termination and cost}
Apply Theorem~\ref{thm:certificate} to every current component, allowing
arbitrary repeated entries and deletions. Issue an accuracy-$\xi$
certificate only when its explicit error and integer thresholds hold
and its defect test passes. On the good geometry and density event,
every current entry is genuine, the certificate is sound and a complete
saturated component passes. The total exceptional probability at epoch
$k$ is below $2\alpha_k$. Summing over $k\ge1$ proves simultaneous
validity with probability at least $1-\delta/4$, leaving room for the
separate launch-rejection budget. This is stronger than the stated
$1-\delta$ guarantee and needs no union bound over stopping rules.

There is a deterministic $k_\xi$ after which the error in
\eqref{eq:density} and the cloud error meet all thresholds. At an epoch
$k\ge k_\xi$ where the raw witness is present and current validation is
good, that witness is rediscovered from fresh clouds and stopping occurs.
Thus
\[
 \{T_\xi>k\}\subset\{K_{\rm raw}>k\}\cup
                       \{\hbox{current validation fails}\}.
\]
Equations \eqref{eq:coupon} and \eqref{eq:spending} prove
\eqref{eq:tail-main}. An erroneous epoch does not add a persistent
constant to this tail. This is the reason for retaining raw proposals
and rebuilding the registry instead of freezing old statistical names.

There are $O(k^3)$ primary accepted preparations at epoch $k$. The free
probability in \eqref{eq:square} gives conditional expected primary
attempts $O(k^3)$, and a negative-binomial bound gives
$C(k^3+\log((k+1)^a/\delta))$ attempts with failure probability
$\delta/[32(k+1)^a]$. Every cloud uses attempted, not accepted, launches.
Since $\nu_k^{-3/2}\le C(k+1)^{3/2}$ for sufficiently large $k$ and
$\log(C/(\nu_k\zeta_{k,j}))\le C\log(C(k+1)/\delta)$, its fixed total
epoch budget is
\begin{equation}\label{eq:epoch-work}
 C(k+1)^{9/2}\log\frac{C(k+1)}{\delta}.
\end{equation}
Summing proves \eqref{eq:cost-main}; the rejection allowances summed
above are less than $\delta/4$. Local squares have fixed radius and their
centers lie within a bounded flight distance of the primary square.
The largest spatial range is consequently the order stated in the main
theorem. Required impact-position tolerance is at most $c\nu_k^3$.

The decision to execute epoch $k$ is measurable before its new randomness.
Multiplying its conditional expected budget \eqref{eq:epoch-work} by
\eqref{eq:tail-main} at $k-1$ yields a summable series when $a>11/2$.
Hence the total expected number of launch attempts before $T_\xi$ is
finite. Measurement bit complexity, machine localization cost and apparatus
travel are not included in that count. If they obey additional polynomial
cost models, their exponents must be added to this summability test.
The theorem asserts neither a minimax rate nor a universal practical
constant. It does replace an unbounded-accuracy boundary/distance oracle
and a continuation-sized fingerprint by counted collision launches.
This completes the proof of Theorem~\ref{thm:main}.

\section{Information, prior results and retained conclusions}\label{sec:comparison}
\subsection{What has changed in the experiment}
The previous local-value theorem used active local value and distance
access: certified Cartesian chart coverage, boundary-height enclosures
at arbitrarily prescribed tolerance, and enclosures for the distance to
\emph{all} solid near a candidate flight. Its effectiveness was relative
to complex-analytic and geometric numerical priors. Its explicit uniform
short-arc fingerprint could be extremely large. Those statements are
retained with their original assumptions; a symbolic finite expression
was not a practical sensor design.

The present experiment requests none of those geometric values. It
uses freely positioned short launches and localized collision coordinates,
and derives confidence geometry from a finite random cloud. Its local
spatial neighborhood is large enough to sample the complete encountered
bodies. Thus this is a change in access and range, not a claim that the
new data are a subset of the old short-patch data. The number of launches
and the required localization tolerance are both displayed. Independent
resettable launches remain stronger than observing one passive trajectory.
The fixed separation and genericity bounds remain priors, not estimates.

The gain in regularity is correspondingly precise. Whole boundary images
are now sampled, so the certificate works with finite $C^{6,\beta}$ bounds
and needs no holomorphic strip, graph disk or continuation modulus. The
class is not merely analytic bodies described with different notation.
For example, at any member with strict margins, sufficiently small
$C^{6,\beta}$ perturbations of its support functions preserve all the
finitely many bounded-copy inequalities and the shape and symmetry gaps.
One may choose a smooth periodic bump, flat at the edge of its support
and nonzero inside, which is not real analytic. These perturbations give
nonanalytic physical members wherever the original strict-margin class
is nonempty. The earlier registered symmetric-table and analytic
short-contact results are not restricted by this new acquisition theorem.

\subsection{Nearest inverse-geometry comparisons}
The canonical relation generated by $W$ has momenta $-W_s,W_t$.
Knowing that relation gives $W$ only up to a constant on its local branch;
an absolute action value and the normalization $A$ are additionally
needed to recover the two densities. Conversely \eqref{eq:action-area}
recovers these from the densities. Differentiation of travel-time data
is a classical inverse-geometry mechanism, not a new abstract principle
claimed here.

Stefanov--Uhlmann--Vasy~\cite{SUV} determine a metric near a strictly
convex boundary point from local boundary distance, modulo diffeomorphism,
and prove lens rigidity under a convex-foliation hypothesis. The unknown
there is a Riemannian metric. Our metric is fixed Euclidean, the unknowns
are periodic reflecting bodies and their translation group, and the
experiment includes internal localized impacts. Neither information set
is silently identified with the other.

Gurfinkel--Noakes--Stoyanov~\cite{GNS23a}, Theorem 1, prove uniqueness
for finite unions of strictly convex obstacles in a containing strictly
convex Riemannian domain under their curvature conditions and equivalence
up to tangency; general position implies the latter condition. Their
subsequent work~\cite{GNS23b}, Theorem 1, removes that tangency-equivalence
hypothesis in dimension at least three under the stated curvature
conditions. Its two-dimensional result is conditional on the admissible-path
property in Theorem 2. These are exterior travelling-time data, triples
of entry point, exit point and length on a containing boundary. We do not
claim their higher-dimensional result without qualifications in the plane,
nor treat our collision-localization apparatus as an exterior travel-time
sensor. Their results do not supply the particular periodic completion
defect or the finite launch accounting studied here; obstacle recovery
from rich scattering information itself is not a novelty claim.

The collision-cloud proof also has a direct stochastic-geometry comparison.
Schneider~\cite{Schneider} studies boundary-sampled random convex hulls in
the plane, including Hausdorff approximation. Prochno--Schuett--Sonnleitner--Werner
\cite{PSSW} give modern Hausdorff-distance results and record the smooth
boundary order $(\log n/n)^2$ in dimension two. We neither claim that
boundary-hull rate nor mollifier interpolation as new. Our use proves
that the requisite boundary coverage is generated by the physical
first-hit experiment without known labels, and combines it with
unknown-solid confidence, period recognition and refreshed stopping.
The finite bound here is deliberately conservative and is not an
asymptotic constant or minimax theorem.

\subsection{Qualification of the retained effective theorem}
To make the earlier numerical-prior convention unambiguous, its effective
recognition theorem uses the following directed inputs: centered support
strip width and upper bound $(\rho,M)$; local graph disk radius and upper
bound $(R,M_g)$; lower curvature $\kappa_-$; upper bridge length $R_+$
and body diameter $D_0$; lower rotation/reflection margin $\eta$, inter-type
shape gap $\sigma$, and physical-body separation $d_0$; and all chart,
time, active and clearance margins required for acquisition. Lower bounds
must be positive and valid, and upper bounds must actually dominate the
corresponding quantities. These analytic inputs belong to that retained
short-arc theorem, not to Definition~\ref{def:prior}.

In the retained value-probe cost statement, the required height error
\emph{tolerance is at most} $c\nu^3$. The ambiguous expression about
accuracy being ``no smaller'' is to be read with this explicit inequality.
The retained per-probe constant contains its fixed value-grid size $N$;
so do its epoch-work constants. Expressions such as $N=2^P$ describe
symbolic computability, not a physically expanded or queried vector.
The new descriptor has size \eqref{eq:fingerprint-size} instead, but
uses measured full-body geometry. Comparing the two sizes without that
sensor distinction would be misleading.

The old order-dependent moment regularization, analytic full-histogram
regularization and present cloud-assisted regularization remain different
experiments. The first uses finitely many types of exact moment germs;
the second measures densities and continues arcs; the third samples full
nearby bodies as well as the normalized endpoint law. No optimality
comparison follows by comparing powers of their different error parameters.

\subsection{Preserved manuscript and scope}
The complete reviewed v25 manuscript is supplied unchanged in
\texttt{retained/v25/}, with every earlier retained manuscript nested
inside it. It includes the analytic effective calibration, local-value
construction, coefficient inverses, count fibers, registered and
unregistered variants, and their original proofs. The original complete
smooth relative-law supplement is also preserved at \texttt{complete/}.
The source map identifies the active main documents; previous build
receipts remain historical and are not evidence of a current execution.
This primary article is self-contained for its new finite-smoothness
certification route. No retained mathematical argument is deleted.

Finally, equality of finite or formal count data, equality of all smooth
Taylor coefficients, and equality of exact analytic count germs are not
interchangeable. Neither the new collision cloud nor its completion
certificate proves a count-only analytic theorem. Nor is the result a
marked-length-spectrum rigidity theorem, a passive-trajectory discovery
theorem, a theorem without positive generic margins, or a claim of
journal acceptance. The new assertions concern the explicitly stated
launch experiment and its reference-free whole-table certificate.

\subsection*{Acknowledgment and responsibility}
AI-assisted analysis contributed to the fixed-aperture inverse,
collision-flux acquisition, finite-smoothness recognition, compensated
support reconstruction, revalidation and source-verification arguments. The named author is responsible for checking the proofs and
references before submission. Finite diagnostics and compilation are
reproducibility evidence, not independent mathematical certification.
\begin{thebibliography}{99}
\raggedright
\bibitem{SUV}
P. Stefanov, G. Uhlmann, and A. Vasy,
\emph{Local and global boundary rigidity and the geodesic X-ray transform
in the normal gauge}, arXiv:1702.03638v3 (2021).
\bibitem{GNS23a}
T. Gurfinkel, L. Noakes, and L. Stoyanov,
\emph{Uniqueness of obstacles in Riemannian manifolds from travelling times},
arXiv:2309.11141v1 (2023).
\bibitem{GNS23b}
T. Gurfinkel, L. Noakes, and L. Stoyanov,
\emph{Rigidity of travelling times for strictly convex obstacles in
Riemannian manifolds}, arXiv:2311.07813v1 (2023).
\bibitem{Schneider}
R. Schneider, \emph{Random approximation of convex sets},
J. Microsc. \textbf{151} (1988), 211--227.
\href{https://doi.org/10.1111/j.1365-2818.1988.tb04682.x}{doi:10.1111/j.1365-2818.1988.tb04682.x}.
\bibitem{PSSW}
J. Prochno, C. Sch\"utt, M. Sonnleitner, and E. M. Werner,
\emph{Random approximation of convex bodies in Hausdorff distance},
Math. Ann. \textbf{392} (2025), 4525--4542.
\href{https://doi.org/10.1007/s00208-025-03186-7}{doi:10.1007/s00208-025-03186-7}.
\end{thebibliography}

\end{document}
